\documentclass[11pt]{article}

\usepackage[a4paper,margin=1in]{geometry}
\usepackage[utf8]{inputenc}
\usepackage[T1]{fontenc}
\usepackage{amsmath,amssymb}
\usepackage{booktabs}
\usepackage{array}
\usepackage{xcolor}
\usepackage{listings}
\usepackage{hyperref}
\usepackage{enumitem}
\usepackage{titlesec}
\usepackage{fancyhdr}

\hypersetup{
  colorlinks=true,
  linkcolor=blue!50!black,
  citecolor=blue!50!black,
  urlcolor=blue!50!black,
  pdftitle={BDH to AttnGrounder: A Mechanistic Study Guide},
  pdfauthor={Sudarshan Sudhakar}
}

\definecolor{codebg}{RGB}{245,245,248}
\definecolor{codekw}{RGB}{0,90,180}
\definecolor{codecomment}{RGB}{110,110,110}

\lstset{
  basicstyle=\ttfamily\small,
  backgroundcolor=\color{codebg},
  commentstyle=\color{codecomment}\itshape,
  keywordstyle=\color{codekw}\bfseries,
  frame=single,
  rulecolor=\color{gray!40},
  framesep=6pt,
  xleftmargin=6pt,
  xrightmargin=6pt,
  breaklines=true,
  showstringspaces=false,
  columns=flexible,
  language=Python,
  morekeywords={ReLU,softmax,sigmoid}
}

\pagestyle{fancy}
\fancyhf{}
\fancyhead[L]{\small BDH $\to$ AttnGrounder}
\fancyhead[R]{\small Study Guide}
\fancyfoot[C]{\thepage}
\renewcommand{\headrulewidth}{0.3pt}

\titleformat{\section}{\Large\bfseries\color{blue!40!black}}{\thesection}{0.5em}{}
\titleformat{\subsection}{\large\bfseries\color{blue!30!black}}{\thesubsection}{0.5em}{}

\newlength{\boxinnerwidth}
\newsavebox{\keyboxbin}
\newsavebox{\qaboxbin}

\newenvironment{keybox}[1][]{%
  \def\tmptitle{#1}%
  \setlength{\boxinnerwidth}{\dimexpr\linewidth-2\fboxsep-2\fboxrule\relax}%
  \begin{lrbox}{\keyboxbin}
  \begin{minipage}{\boxinnerwidth}
  \ifx\tmptitle\empty\else\noindent\textbf{\tmptitle}\par\smallskip\fi
}{%
  \end{minipage}
  \end{lrbox}%
  \par\addvspace{8pt}\noindent\fcolorbox{blue!35!black}{blue!4}{\usebox{\keyboxbin}}\par\addvspace{8pt}%
}

\newenvironment{qabox}[1]{%
  \setlength{\boxinnerwidth}{\dimexpr\linewidth-2\fboxsep-2\fboxrule\relax}%
  \begin{lrbox}{\qaboxbin}
  \begin{minipage}{\boxinnerwidth}
  \noindent\textbf{#1}\par\smallskip
}{%
  \end{minipage}
  \end{lrbox}%
  \par\addvspace{8pt}\noindent\fcolorbox{orange!60!black}{yellow!6}{\usebox{\qaboxbin}}\par\addvspace{8pt}%
}

\newcommand{\R}{\mathbb{R}}

\title{\textbf{BDH $\to$ AttnGrounder}\\[4pt]
\large A Mechanistic Study Guide: What We Built and the Math Behind It}
\author{Sudarshan Sudhakar}
\date{\today}

\begin{document}
\maketitle

\begin{keybox}[{One-sentence pitch}]
AttnGrounder originally grounds a phrase by having every image region compute a
softmax-weighted average over the command's word vectors. We replaced that one
attention module with BDH's mechanism --- lift regions and words into a
high-dimensional sparse space, build one explicit associative-memory matrix from
the words via an outer product, let every region read from that same memory, and
combine what a region already represents with what it retrieved \emph{multiplicatively}
rather than by averaging --- to test whether this sharper, memory-based readout
disambiguates scenes with multiple same-class objects better than softmax does.
\end{keybox}

\tableofcontents

% =====================================================================
\section{The Baseline: AttnGrounder's Original Attention}

File: \texttt{external/AttnGrounder/model/grounding\_model.py}, method
\texttt{text\_attn} ($\sim$line 199).

\paragraph{Inputs / outputs.}
\begin{itemize}[leftmargin=1.4em]
  \item \texttt{image\_feat}: $(B, C, H, W)$ --- a grid of visual region features,
        one $C$-dim vector per spatial cell ($C=512$).
  \item \texttt{lang\_feat}: $(B, T, C)$ --- one feature vector per word, from a
        BiLSTM over GloVe embeddings.
  \item Output 1, \texttt{lang\_feat\_attn}: $(B, C, H, W)$ --- for every region,
        ``what does the text say about this region,'' same shape as
        \texttt{image\_feat} so it can be concatenated back with it.
  \item Output 2, \texttt{beta}: $(B, H, W)$ --- a heatmap in $[0,1]$, one scalar
        per region, supervised directly by a BCE loss against the ground-truth box
        mask (an auxiliary loss forcing the attention to light up on the right region).
\end{itemize}

\paragraph{The math.} Let $R = H \cdot W$ be the number of regions and $T$ the number
of words. With $G \in \R^{R\times C}$ the flattened region features and
$Q \in \R^{T \times C}$ the word features:
\begin{align}
  S &= G\,Q^{\top} && \in \R^{R \times T} \quad \text{(raw region--word dot products)}\\
  A &= \operatorname{softmax}_{T}(S) && \in \R^{R \times T} \quad \text{(normalize over words, per region)}\\
  \text{lang\_feat\_attn} &= A\,Q && \in \R^{R \times C} \quad \text{(weighted average of word vectors)}\\
  \beta &= \sigma\!\Big(\textstyle\sum_{t} S_{\cdot,t}\Big) && \in \R^{R} \quad \text{(``how strongly does this region match \emph{any} word'')}
\end{align}

This is standard \textbf{cross-attention}: every region asks ``which words describe
me?'', gets a softmax distribution over the $T$ words, and takes a weighted average
of their embeddings. It has \textbf{zero learned parameters} --- purely
dot-product $\to$ softmax $\to$ sigmoid on features already produced upstream by the
CNN and BiLSTM.

\begin{keybox}[Why this is a weak spot]
Softmax attention gives every region a smooth, \emph{convex combination} of word
vectors. If the phrase is ``the car on the left'' and there are two cars in the
scene, both car-regions score similarly against the word ``car'' --- softmax just
averages; it has no mechanism to sharply separate ``this car'' from ``that car''
beyond whatever the raw upstream features already encode. There is no notion of a
persistent, content-addressable memory --- it is a single, one-shot, low-capacity
read.
\end{keybox}

% =====================================================================
\section{What BDH Is: Intuition, Then Math}

BDH = ``Baby Dragon Hatchling'' (ground-truth code: \texttt{external/bdh/bdh.py};
paper \texttt{papers/2509.26507v1.md}). Mechanically, forget the biological framing
for defense purposes: it is a \textbf{linear-attention Transformer with an explicit
outer-product associative memory}, built from three ideas.

\subsection{Lift into a big, sparse, positive space}

Instead of attending directly in the model's working dimension $d$ (e.g.\ 512), BDH
first \textbf{lifts} every feature vector into a much higher dimension
$n = \texttt{mult} \times d$ (we use $\texttt{mult}=4 \Rightarrow n=2048$) via a
learned linear map $E \in \R^{d\times n}$, followed by ReLU:
\begin{equation}
  x_{\text{sparse}} = \operatorname{ReLU}(x\,E) \in \R^{n}, \qquad n \gg d .
\end{equation}

\begin{keybox}[Intuition]
In a low dimension $d$, two similar-but-different concepts (car \#1, car \#2) sit
close together and small distinguishing signals get lost. Projecting into a much
higher dimension and clipping negatives with ReLU is exactly what sparse coding /
kernel feature maps do: it is far easier to make two points \emph{linearly
separable} and \emph{sparse} (few active neurons each) in a big space than a small
one. Only a handful of the $n$ neurons fire for any given input, and different
inputs tend to activate mostly non-overlapping subsets of neurons --- like a hash.
(BDH's paper motivates this via a Hebbian/neuronal analogy --- ``neurons that fire
together wire together'' --- but functionally it is a sparse, positive, overcomplete
ReLU feature map, in the same family as ReLU kernel feature maps used in linear
attention.)
\end{keybox}

\subsection{Attention becomes an explicit associative memory}

Standard attention computes $\operatorname{softmax}(QK^{\top})V$: normalize scores
over keys, then average values. That is a \emph{soft lookup table with no
persistent state} --- every query re-derives its weights from scratch against all
keys.

BDH instead builds a memory matrix explicitly. With $K_{\text{sparse}}, Q_{\text{sparse}}
\in \R^{\cdot \times n}$ the lifted keys/queries and $V \in \R^{\cdot \times d}$ the
values:
\begin{align}
  \rho &= K_{\text{sparse}}^{\top} V && \in \R^{n \times d} \quad \text{(outer-product accumulation over \emph{all} keys/values)} \\
  \text{read} &= Q_{\text{sparse}}\,\rho && \in \R^{\cdot \times d} \quad \text{(each query retrieves by dotting into that memory)}
\end{align}

This is \textbf{linear attention}: no softmax, just a weighted sum of outer products
$\sum_i k_i \otimes v_i$, then each query multiplies against that sum. Algebraically
it is equivalent to $Q_{\text{sparse}}K_{\text{sparse}}^{\top}V$ (re-associated for
efficiency) --- i.e.\ ``unnormalized linear attention.'' What matters for the
grounding use case is the \textbf{framing}: $\rho$ is a literal memory matrix that
can be \emph{built once} (accumulate keys and values into it) and \emph{read many
times} (any query dots into it) --- content-addressable storage, like a Hopfield
network or a linear key--value store, rather than a fresh softmax computed per
query.

\begin{keybox}[Why this maps well onto grounding]
In our use, the ``keys/values'' are the words of the command and the ``queries''
are image regions. Building the memory once from all words, then letting every
region address into it, is literally what a \emph{referring expression} should
behave like: a small declarative store (``car'', ``left'', ``white'') that every
candidate region queries against.
\end{keybox}

\subsection{A multiplicative gate, not a linear blend}

After the memory read, BDH does not simply decode \texttt{read} back to $d$
dimensions. It re-lifts the read result into the same sparse $n$-dim space and
\textbf{multiplies it, element-wise}, against the original query's own sparse
address:
\begin{align}
  \text{read}_{\text{sparse}} &= \operatorname{ReLU}(\text{read}\,E_v) && \in \R^{\cdot \times n} \\
  \text{gated} &= Q_{\text{sparse}} \odot \text{read}_{\text{sparse}} && \in \R^{\cdot \times n} \quad \text{(elementwise product)} \\
  \text{out} &= \text{gated}\,D && \in \R^{\cdot \times d} \quad \text{(decode)}
\end{align}

\begin{keybox}[Why multiplicative matters]
A weighted \emph{sum} (softmax attention) can only ever interpolate between
existing value vectors --- it is linear in the values. A gated \emph{product}
between ``what I was looking for'' ($Q_{\text{sparse}}$) and ``what I found''
($\text{read}_{\text{sparse}}$) creates a \textbf{bilinear} interaction: the output
carries signal only where a neuron is active in \emph{both} the query's own
representation and the retrieved memory. This is a sharper, AND-like combination
(vs.\ softmax's OR-like/average combination) and is the mechanism most directly
responsible for BDH's claimed ability to keep near-duplicate concepts (e.g.\ two
same-class objects) separable rather than blurred together.
\end{keybox}

\paragraph{The whole engine, in one line:} \emph{lift $\to$ build memory
(outer product) $\to$ read $\to$ gate multiplicatively $\to$ decode.} Every layer of
the original BDH language model (\texttt{external/bdh/bdh.py}, class \texttt{BDH})
repeats this, plus RoPE positional encoding and a causal mask (it is a next-token
predictor). Our job was to strip the parts that only make sense for autoregressive
text and repurpose the memory mechanism for \emph{cross-modal grounding}, which has
no sequence order between ``regions'' and ``words'' at all.

% =====================================================================
\section{What We Built: \texttt{bdh\_grounding/bdh\_fusion.py}}

We wrote \texttt{BDHVisualTextAttention}, a \textbf{drop-in replacement} for
\texttt{text\_attn} --- identical inputs/outputs --- selected in
\texttt{grounding\_model.py} via \texttt{variant == "bdh"}
(\texttt{self.bdh\_attn = BDHVisualTextAttention(\ldots)}, one instance shared
across AttnGrounder's 3 FPN scales).

Three variants exist in code (the \texttt{mode} flag), but \textbf{only mode
``C'' is validated and used for reported results} --- A and B are ablations kept
for future comparison.

\subsection{Mode C (primary, validated): symmetric, non-causal}

This is the direct translation of Section~2 into cross-attention. Let
$G \in \R^{R\times d}$ be the region features and $Q \in \R^{T\times d}$ the word
features (both LayerNorm'd first):
\begin{align}
  G_q &= \operatorname{ReLU}(G\,E) && \in \R^{R \times n} \quad \text{region \emph{query} addresses} \\
  K_k &= \operatorname{ReLU}(Q\,E) && \in \R^{T \times n} \quad \text{word \emph{key} addresses} \\
  \rho &= K_k^{\top} Q && \in \R^{n \times d} \quad \text{memory, built once from \emph{all} $T$ words} \\
  \text{read} &= (G_q\,\rho) \cdot \tfrac{1}{\sqrt n} && \in \R^{R \times d} \quad \text{each region reads} \\
  \text{read}_{\text{sparse}} &= \operatorname{ReLU}\big(\operatorname{LN}(\text{read})\,E_v\big) \\
  \text{gated} &= G_q \odot \text{read}_{\text{sparse}} && \in \R^{R \times n} \quad \text{multiplicative gate} \\
  T' &= \operatorname{LN}(\text{gated}\,D) && \in \R^{R \times d} \;\to\; (B,C,H,W) \\
  \beta &= \sigma\big(\operatorname{Linear}_{d\to1}(\text{read})\big) && \in \R^{R} \;\to\; (B,H,W)
\end{align}

\begin{lstlisting}[caption={Mode C forward pass (\texttt{bdh\_grounding/bdh\_fusion.py:\_forward\_C}), abbreviated}]
Gq  = relu(G @ E)                 # (B, R, n) region query addresses
Kk  = relu(Q @ E)                 # (B, T, n) word key addresses
rho = Kk.transpose(1, 2) @ Q      # (B, n, d) associative memory, built from ALL words
read = (Gq @ rho) * scale         # (B, R, d) each region reads from memory
read_sparse = relu(ln(read) @ Ev)
gated = Gq * read_sparse          # (B, R, n) multiplicative gate
Tprime = ln(gated @ Dx)           # (B, R, d) -> reshaped to (B, C, H, W)
beta = sigmoid(beta_head(read))   # learned scalar head -> (B, H, W)
\end{lstlisting}

\paragraph{Mapping onto Section 2/3.}
\begin{itemize}[leftmargin=1.4em]
  \item \textbf{Regions are queries, words are keys/values.} This is the
        \emph{opposite} direction from BDH's original language-model use (where a
        token queries \emph{past} tokens), but here there is no ``past'':
        grounding is symmetric / non-causal, so we drop BDH's causal mask and RoPE
        entirely --- word order within the sentence doesn't gate what a region is
        allowed to attend to.
  \item $\rho = K^{\top}Q$ is built \textbf{once per forward pass} from all $T$
        words --- literally ``memorize the whole command'' --- then every one of
        the $R = H{\times}W$ regions reads from that \emph{same} memory
        independently. This replaces the baseline's per-region softmax-over-words
        with per-region dot-product-into-memory.
  \item The gate $G_q \odot \text{read}_{\text{sparse}}$ is the bilinear AND from
        Section~2.3: a region's own sparse signature must overlap with what its
        memory read up in order to produce a large output, rather than the
        baseline's linear softmax average.
\end{itemize}

\begin{qabox}[Why did \texorpdfstring{$\beta$}{beta} need a learned head?]
Unlike the baseline's parameter-free $\sigma(\sum S)$, mode C uses
\texttt{sigmoid(Linear($d\to1$)(read))}. BDH's ReLU-sparse dot products
($G_q$, $K_k$) are all $\geq 0$ by construction, so a raw sum of them can never dip
meaningfully below the sigmoid's midpoint. The auxiliary BCE mask loss needs
$\beta \to 0$ on background regions --- unreachable without a learned bias/weight to
shift it. This is a $\sim(d{+}1)$-parameter head, negligible in size but necessary:
BDH is \emph{not} fully parameter-free here, by design, for exactly this reason.
\end{qabox}

\subsection{Fusion trim (parameter budget)}

Baseline AttnGrounder concatenates three streams before its output convolution:
$[\,\text{visual},\ \text{lang\_feat\_attn},\ \beta \odot \text{visual}\,]$, i.e.\
$\texttt{embin\_size} = \texttt{emb\_size} \times 3$. For the BDH variant we fuse
only $[\,\text{visual},\ \text{lang\_feat\_attn}\,]$ (drop the $\beta\odot$visual
product stream): $\texttt{embin\_size} = \texttt{emb\_size} \times 2$. This is a
parameter-budget decision: it guarantees the BDH-variant model's total parameter
count does not exceed the 75.84M baseline, so any accuracy difference is not
confounded by ``we just gave the model more capacity.'' $\texttt{mult}=2$ is the
default operating point, balancing BDH's need for $n \gg d$ (expansion gives sparse
separability, but costs $n{\times}d$ params per matrix) against that budget.

\subsection{Modes A and B: implemented ablations}

\begin{description}[leftmargin=1.6em, style=nextline]
  \item[Mode A] Text is first passed through a \emph{causal} BDH self-attention
    pass over itself (words attend only to earlier words), producing a
    contextualized memory, and \emph{then} regions query that. Tests whether
    giving the text module its own BDH ``understanding pass'' before grounding
    helps --- an in-context-retrieval framing closer to BDH's original LM use.
  \item[Mode B] The most literal port of the original \texttt{bdh.py} forward
    pass: concatenate $[\text{words} \,;\, \text{regions}]$ into \emph{one}
    sequence, run genuine causal BDH self-attention (with RoPE) over the whole
    thing, then read out only the region positions at the end --- ``vanilla BDH''
    with regions appended as extra, causally-later tokens.
\end{description}

Mode C was chosen as primary because grounding has no real notion of sequence
order between ``words'' and ``regions'': forcing a causal mask (A, B) throws away
information for no modeling reason, whereas C's symmetric non-causal memory
matches the task structure directly.

\subsection{Cross-scale memory (``growing'' prior)}

AttnGrounder runs \texttt{text\_attn} \textbf{independently at 3 FPN scales}
(different $H\times W$ resolutions). We added an optional \texttt{region\_prior}
argument: each scale's own \texttt{lang\_feat\_attn} output can be fed forward
(upsampled + added as a residual) into the next, finer scale's region features
before that scale builds its own memory --- the ``growing memory across scales''
idea, orthogonal to the A/B/C mode choice. Off by default
(\texttt{region\_prior=None} reproduces the original per-scale-independent
behavior exactly), toggled via \texttt{bdh\_growing\_scales}.

% =====================================================================
\section{Side-by-Side: What Actually Changed, Mathematically}

\begin{center}
\renewcommand{\arraystretch}{1.35}
\begin{tabular}{p{2.6cm}p{5.5cm}p{6.3cm}}
\toprule
 & \textbf{Baseline (softmax cross-attn)} & \textbf{BDH (mode C)} \\
\midrule
Score function & $QK^{\top}$ (dot product in raw $d$-dim space) & $Q_{\text{sparse}}(K_{\text{sparse}}^{\top}V)$ --- score pre-aggregated into a memory matrix \\
Normalization & softmax over words (sums to 1) & none --- unnormalized linear attention, scaled by $1/\sqrt n$ \\
Combine rule & weighted \textbf{sum} of value vectors (linear, convex) & \textbf{outer-product memory read}, then elementwise-\textbf{multiplied} with the query's own sparse code (bilinear) \\
Feature space & raw $d$-dim (e.g.\ 512) & lifted $n$-dim sparse ReLU space ($n = \texttt{mult}\cdot d$, e.g.\ 2048), $n \gg d$ \\
Params & 0 & $E, E_v, D$ ($\approx 3nd$) $+$ tiny $\beta$ head \\
Complexity & $O(R\,T\,d)$ & $O(Rn) + O(Tn) + O(nd)$ --- memory $\rho$ built once, reused by every region \\
\bottomrule
\end{tabular}
\end{center}

\begin{keybox}[The one sentence version]
Softmax attention computes a fresh weighted average per region; BDH compiles the
entire command into one small memory matrix once, then lets every region address
into it through a sparse, multiplicative (AND-like) readout instead of a smooth
average.
\end{keybox}

% =====================================================================
\section{The Hypothesis: Why This Should Help Ambiguous Scenes}

This is the hypothesis actually being tested. Same-class ambiguity (two cars in
frame, command says ``the car on the left'') is exactly where \textbf{softmax
averaging degrades}: if both car-regions score similarly against the word
``car'', softmax spreads a similar-shaped distribution across both, and the linear
weighted-sum blends their word evidence almost identically.

BDH's sparse lift makes it more likely that subtly different regional features
(car-on-left vs.\ car-on-right, differing e.g.\ in local context) activate
\textbf{different subsets} of the $n$ sparse neurons, and the multiplicative gate
then produces sharply different outputs for regions whose sparse codes do not
overlap the retrieved memory the same way --- i.e.\ BDH has structurally more room
to \emph{disambiguate} than a low-dimensional softmax does.

This is why evaluation is \textbf{stratified}: \texttt{analysis/stratified\_eval.py}
splits AP50 into ``ambiguous'' ($\geq 1$ other same-class object in the scene, from
nuScenes metadata) vs.\ ``unambiguous,'' specifically to check whether any BDH gain
\emph{concentrates} where the hypothesis predicts it should, rather than being a
uniform, unexplained bump.

% =====================================================================
\section{Anticipated Questions}

\begin{qabox}[Is this the actual BDH paper's architecture?]
No --- it is BDH's \emph{core mechanism} (sparse lift + outer-product memory +
multiplicative gate) repurposed as a cross-attention module. We dropped what only
applies to autoregressive language modeling: the causal mask, RoPE positional
encoding, and the token embedding / LM head. Mode B is the closest to ``literal
BDH''; mode C is the adapted version and the one we report.
\end{qabox}

\begin{qabox}[Isn't \texorpdfstring{$\rho = K^{\top}Q$}{rho = KTQ} just a matmul reordering trick?]
Algebraically $Q(K^{\top}V) = (QK^{\top})V$, so yes --- it is the identity linear
attention always uses. The ``memory'' framing matters because $\rho$ is a fixed
$(n,d)$ object built \emph{once} per forward pass and reused by every one of the
$R$ region queries --- it does not recompute anything per-region, unlike softmax's
$QK^{\top}$, which is inherently a full $R \times T$ pairwise matrix. It is a
legitimate architectural distinction (persistent, content-addressable state vs.\
per-query recomputation), not just terminology.
\end{qabox}

\begin{qabox}[What's the parameter budget vs.\ the baseline?]
Kept $\leq 75.84$M via the fusion trim (Section 3.2); \texttt{mult} and
\texttt{share\_qv\_encoder} are the knobs that trade memory capacity ($n$) against
parameter count. \texttt{share\_qv\_encoder=true} ties the query-address and
value-path encoders ($E = E_v$), reclaiming $n \times d$ parameters.
\end{qabox}

\begin{qabox}[Did you validate modes A and B too?]
Implemented, not the headline result --- kept as future ablation work; only mode
C has been trained and evaluated end-to-end.
\end{qabox}

\begin{qabox}[What about TransVG?]
A second, independent integration
(\texttt{bdh\_grounding/bdh\_selfattn.py}) swaps BDH's mechanism into TransVG's
transformer \emph{self}-attention (not cross-attention) --- same
lift$\to$memory$\to$gate recipe, but built signature-compatible with
\texttt{nn.MultiheadAttention} so it drops straight into TransVG's encoder layers,
with padded tokens masked out of the memory ($\rho$) so they never contribute.
This was the cross-architecture validation experiment: does the effect show up in
a second, unrelated grounding architecture, not just AttnGrounder.
\end{qabox}

% =====================================================================
\section{Summary Statement}

\begin{keybox}
``AttnGrounder originally grounds a phrase by having every image region do a
softmax-weighted average over the command's word vectors --- a smooth,
parameter-free lookup. We replaced that one module with BDH's mechanism: lift both
regions and words into a much higher-dimensional sparse space, build one explicit
associative-memory matrix from all the words via an outer product, let every
region read from that same memory, then combine the region's own sparse code with
what it retrieved \emph{multiplicatively} rather than by averaging. It's
mathematically linear attention with a sparse ReLU feature map and a bilinear gate
instead of softmax's convex combination, and the hypothesis is that this sharper,
AND-like combination rule should specifically help disambiguate scenes with
multiple same-class objects --- which is why we stratify evaluation by scene
ambiguity rather than just reporting one aggregate AP50.''
\end{keybox}

\vspace{1em}
\noindent\textbf{Where to look yourself, to re-derive any line:}
\begin{itemize}[leftmargin=1.4em, itemsep=2pt]
  \item Baseline attention: \texttt{external/AttnGrounder/model/grounding\_model.py:199}
  \item BDH ground truth: \texttt{external/bdh/bdh.py} (class \texttt{BDH}, \texttt{forward})
  \item AttnGrounder integration: \texttt{bdh\_grounding/bdh\_fusion.py} (\texttt{\_forward\_C} is the one that matters)
  \item TransVG integration: \texttt{bdh\_grounding/bdh\_selfattn.py}
  \item Design rationale: project memory notes on the locked module design and evaluation plan
\end{itemize}

\end{document}
